\documentclass[a4paper]{article}
\usepackage[pdftex]{graphicx}
\usepackage[utf8]{inputenc}
\usepackage{enumerate}
\usepackage{siunitx}
\sisetup{locale=DE} 
\usepackage{href-ul}
\hypersetup{
	colorlinks=true,
	linkcolor=blue,
	urlcolor=blue}
\usepackage{icomma}
\usepackage{amssymb}
\usepackage{geometry}
\geometry{a4paper, top=15mm, left=15mm, right=15mm, bottom=15mm,
	headsep=10mm, footskip=12mm}

\begin{document}
{\bf Kreisbewegung}
\begin{enumerate}[1.]

%Aufgabe 1
{\bf \item Eine Kugel wird an einer Schnur der Länge $l=\SI{40}{\centi\meter}$ befestigt und mit einer Frequenz 
$f=\SI{1,5}{\hertz}$ herumgeschleudert.} Dabei beschreibt sie eine horizontale Kreisbahn mit konstanter Winkelgeschwindigkeit $\omega$. 

\begin{enumerate}[a)]
\item Berechne Winkel-- und Bahngeschwindigkeit der Kugel.
\item Als die Schnur reißt, wird die Kugel waagrecht weggeschleudert. In welcher Entfernung und unter welchem Winkel trifft sie auf dem Boden auf, wenn sie sich zum Zeitpunkt des Reißens in einer Höhe von $h=\SI{1,7}{\meter}$ befand?
\end{enumerate}

{\bf \item Lunerald rast mit seinem E-Bike  über einen Feldweg.} Sein Tachometer  zeigt $\SI{36}{\kilo\meter\over \hour}$ an. Mit welcher Frequenz drehen sich die Räder, wenn der Raddurchmesser 28 Zoll beträgt ? 
[1 Zoll $\approx \SI{2,5}{\centi\meter}$].

{\bf \item Eine Audio-CD enthält einen bespielbaren Kreisring} mit einem Außenradius von $r_a=\SI{5,8}{\centi\meter}$ und einem Innenradius von $r_i=\SI{2,2}{\centi\meter}$ Sie wird von innen nach außen abgespielt, dabei ist die Geschwindigkeit des Lesekopfs über der Scheibe immer konstant. Die Drehzahl f ist zu Beginn $\SI{500}{1\over\minute}$. 
\begin{enumerate}[a)]
	\item Wie groß ist die Drehzahl am Ende, wenn fast die gesamte CD abgespielt ist?
	\item Schätze ab, wie ,,lang" (in Metern) die abgelesene Spur der gesamten CD ist, bei 74 Minuten Abspieldauer. 
	\item Wie breit ist diese Spur dann? 
\end{enumerate}


{\bf \item Der Drehzahlmesser von Luneralds Sportwagen} zeigt bei einer Fahrt auf der Autobahn  $\SI{4500}{1\over\minute}$. Das Getriebe übersetzt 3 zu 1. Wie schnell fährt Lunerald, wenn der Reifendurchmesser seines Wagens 60 cm beträgt?

{\bf \item Der Minutenzeiger der Big Ben-Turmuhr } hat eine Länge von 3 m. Welche Strecke überstreicht er in einer Minute / einer Sekunde ?

{\bf \item Mit welcher Bahngeschwindigkeit} bewegt sich die Mündung des Amazonas (Äquator) um den Erdmittelpunkt? [$R_{Aeq}=\SI{6378}{\kilo\meter}$]

{\bf \item In einem Wasserstoffatom umkreist das Elektron} das Proton mit einer Bahngeschwindigkeit von\\ $v=\SI{2,2e6}{\meter\over \second}$. Pro Sekunde finden $\SI{6,6e15}{}$ Umrundungen statt. Welchen Radius kann man für dieses Atom annehmen?

{\bf \item Planeten }

\renewcommand{\arraystretch}{2}
\begin{tabular}{|c|c|c|c|c|c|c|c|c|}
\hline
Körper & Merkur & Venus & Erde & Mars & Jupiter & Saturn & Uranus & Neptun \\
\hline
mittlerer Sonnenabstand in $10^6$ ~km & 57,9 & & 149,6 & 227,9 & & 1427 & 2871 & \\
\hline
Umlaufzeit in a &  & 0,615 & 1,000 & &11,86 &29,45 & & 164,8\\
\hline
mittlere Bahngeschwindigkeit in $\SI{}{\kilo\meter\over \second}$ & 47,9 & 35,1 &  & 24,1 & 13,1&  & 6,80 & 5,43\\
\hline
\end{tabular}

{\bf \item Mit der Foucault'schen Drehspiegelmethode lässt sich die Lichtgeschwindigkeit $c$ messen:}
  
\noindent
\begin{minipage}{0.65\textwidth}
	 Ein Laserstrahl wird auf einen mit $\omega$ schnell rotierenden Drehspiegel gerichtet. An einem weiteren Spiegel in einer Entfernung $d$ wird der Strahl zurück reflektiert und erreicht nach einer Laufzeit $\Delta t = {2d \over c}$ wieder den Drehspiegel. Dieser hat sich zwischenzeitlich um einen Winkel $\alpha=\omega \cdot \Delta t$ weiter gedreht. Dadurch bildet sich ein Winkel $2\alpha$ zwischen dem einfallenden und dem ausfallenden Lichtstrahl, der gemessen werden kann.
	 \begin{enumerate}[a)]
	 	\item  In einem Experiment  beträgt der beobachtete Winkel $2\alpha=\SI{0,0040}{\degree}$.
	 	Um welchen Winkel im Bogenmaß hat sich der Drehspiegel weitergedreht?
	 	\item Wie groß errechnet sich hieraus dann die Lichtgeschwindigkeit $c$, wenn $d=\SI{8,0}{\meter}$ ist und der Drehspiegel  mit einer Frequenz von $f=\SI{100}{\hertz}$ rotiert?
	 \end{enumerate}
\end{minipage}
\hfill
\begin{minipage}{0.35\textwidth}
	\includegraphics[width=6 cm]{foucault414}
\end{minipage}

\end{enumerate}
\fbox{
	\begin{minipage}{0.5\textwidth}
		Zur Lösung bitte \href{https://www.okuyakl.de/phys/p10kreaL414/ll414.pdf}{hier klicken} oder den QR-Code scannen.\\
	Weitere Arbeitsblätter gibt es unter 
	
	\href{https://www.okuyakl.de}{www.okuyakl.de}
	\end{minipage}
	\hfill
	\begin{minipage}{0.4\textwidth}
		\includegraphics[width=1.5 cm]{../../viecher/zwe03}
		\includegraphics[width=3 cm]{qr414}
		\includegraphics[width=2 cm]{../../viecher/afanticon1}
		
	\end{minipage}}

\end{document}